\section{Introduction}
The ISAR kernel provides a unified, substrate-independent framework that bridges discrete combinatory logic (symbolic rewriting) and continuous analysis (neural networks and physical systems). In this blueprint, we formalize the continuous universal approximation and representation space of the ISAR kernel, showing that all continuous functions $\mathbb{R}^d \to \mathbb{R}^k$ are uniquely represented as limit addresses in the completion of the quotient address space.

\section{Continuous Morphism and Address Space}

We define the basic structures of the continuous ISAR analysis:

\begin{definition}[Activation]
An activation function $\sigma$ is a continuous real-valued function:
$$\sigma \in C(\mathbb{R}, \mathbb{R})$$
\end{definition}

\begin{definition}[Non-Polynomial Activation]
An activation function is non-polynomial if it cannot be represented as a polynomial with real coefficients:
$$\forall \text{coeffs} \in \text{List } \mathbb{R}, \quad \sigma \neq \text{evalPoly coeffs}$$
\end{definition}

\begin{definition}[RawAddress]
A raw parameter configuration space $\text{RawAddress } d\ k$ is a tuple representing a neural network configuration:
$$\text{RawAddress } d\ k := \Sigma (N, T : \mathbb{N}), \quad (\text{Fin } T \to \text{Fin } 4 \to \mathbb{R}) \times C(\mathbb{R}^d, \text{GridState } N) \times C(\text{GridState } N, \mathbb{R}^k)$$
where $\text{GridState } N := \text{EuclideanSpace } \mathbb{R} (\text{Fin } N \times \text{Fin } 4)$.
\end{definition}

\begin{definition}[AddressEq]
Two raw addresses $\theta_1, \theta_2 \in \text{RawAddress } d\ k$ are observationally/functionally equivalent ($\theta_1 \sim \theta_2$) if they realize the same continuous function:
$$\text{realizeRaw } d\ k\ \sigma\ \theta_1 = \text{realizeRaw } d\ k\ \sigma\ \theta_2$$
\end{definition}

\begin{definition}[KernelAddress]
The quotient parameter space modulo observational equivalence:
$$\text{KernelAddress } d\ k\ \sigma := \text{Quotient } (\text{addressSetoid } d\ k\ \sigma)$$
\end{definition}

\section{Universal Approximation}

We state the core density property of the ISAR update on compact domains:

\begin{axiom}[ISAR\_UAT]
\lean{ISAR.ISAR_UAT}
\uses{Activation.nonPolynomial, RawAddress}
Let $K \subset \mathbb{R}^d$ be a compact set, $f \in C(\mathbb{R}^d, \mathbb{R}^k)$ a target continuous function, $\sigma$ a non-polynomial activation, and $\varepsilon > 0$. Then there exists a raw address $\theta \in \text{RawAddress } d\ k$ such that:
$$\forall x \in K, \quad \| \text{realizeRaw } d\ k\ \sigma\ \theta\ x - f(x) \| < \varepsilon$$
\end{axiom}

\section{Dense Inclusion and Limit Completion}

\begin{axiom}[KernelAddressLimit \& continuousRealizationLimit]
\lean{ISAR.KernelAddressLimit, ISAR.continuousRealizationLimit}
We axiomatize the metric completion of `KernelAddress` as `KernelAddressLimit` and the continuous realization map on the completion.
\end{axiom}

\begin{axiom}[kernelAddressEmbedding \& continuousRealizationLimit\_coe]
\lean{ISAR.kernelAddressEmbedding, ISAR.continuousRealizationLimit_coe}
The canonical inclusion map $i : \text{KernelAddress} \to \text{KernelAddressLimit}$ and the commuting property $\hat{f}(i(q)) = f(q)$.
\end{axiom}

\begin{theorem}[kernelAddressEmbedding\_injective]
\lean{ISAR.kernelAddressEmbedding_injective}
The completion embedding is injective:
$$\forall q_1, q_2 \in \text{KernelAddress}, \quad i(q_1) = i(q_2) \implies q_1 = q_2$$
\end{theorem}

\begin{theorem}[kernelAddressEmbedding\_dense]
\lean{ISAR.kernelAddressEmbedding_dense}
The completion embedding has dense range (with respect to uniform convergence on compact domains).
\end{theorem}

\begin{axiom}[topological\_extension\_bijection]
\lean{ISAR.topological_extension_bijection}
An injective map with a dense range into a complete metric space extends uniquely to a bijection on the completion of its domain.
\end{axiom}

\begin{theorem}[ISAR\_representation]
\lean{ISAR.ISAR_representation}
\uses{topological_extension_bijection, kernelAddressEmbedding_injective, kernelAddressEmbedding_dense}
Every continuous function $f \in C(\mathbb{R}^d, \mathbb{R}^k)$ has a unique address $\theta_{\text{limit}} \in \text{KernelAddressLimit}$ such that:
$$\text{continuousRealizationLimit } d\ k\ \sigma\ \theta_{\text{limit}} = f$$
\end{theorem}

\section{Physical System Universality}

\begin{theorem}[physical\_system\_has\_address]
\lean{ISAR.physical_system_has_address}
\uses{ISAR_representation}
Every physical system represented as an admissible continuous kernel has a unique address in the continuous completion limit space.
\end{theorem}

\subsection{Physical and Statistical Universality Synthesis}
The mathematical architecture establishes four fundamental system correspondences:
\begin{enumerate}
    \item \textbf{Fisher-Suffizienz}: When $K$ represents a stochastically modeled system, the equivalence class in `InvariantLayer` modulo observational equivalence acts as the \textbf{sufficient statistic} in the sense of Fisher, capturing all observable information.
    \item \textbf{Turing-Stabilität (Gray-Scott)}: The nilpotent matrix update ($K^2 = 0$) limits continuous updates to first-order flows: $\exp(\varepsilon K) = I + \varepsilon K$, preventing chaotic runaway and stabilizing dissipative Turing patterns.
    \item \textbf{Ising-Kopplung}: The block-diagonal parameter matrix with adjacency $A$ models nearest-neighbor spin coupling, where thermodynamic criticality represents the boundary between stable nilpotent and invertible dynamics.
    \item \textbf{Lorenz-Attraktor}: Chaotic trajectories embed naturally in $4$-dimensional spaces (`Fin 4`), where self-similar dynamics correspond to fixed points of the realization map under temporal rescalings.
\end{enumerate}
